% ORGANIZACIÓN DE COMPUTADORAS
% SEGUNDO CUATRIMESTRE 2017
% PRIMER EXAMEN PARCIAL

\documentclass[12pt,a4paper]{article}
\usepackage[spanish]{babel}
\selectlanguage{spanish}
\usepackage[utf8]{inputenc}

\usepackage{graphicx}
\usepackage{multirow}
\usepackage{xcolor}
\usepackage{fancyhdr}
\renewcommand{\headrulewidth}{0pt}

\setlength{\textwidth}{170mm}
\setlength{\oddsidemargin}{-5.4mm}
\setlength{\evensidemargin}{-5.4mm}
\setlength{\textheight}{247mm}
\setlength{\topmargin}{-15.4mm}
\setlength{\headheight}{11mm}
\setlength{\headsep}{2mm}

% CARPETA CON IMAGENES DE ENCABEZADO
\graphicspath{{estilo/}}

% DEFINICIÓN DE COMANDOS

\newcommand{\setyear}[1]{\year=#1}
\newcommand{\universidad}{Universidad Nacional del Sur}
\newcommand{\departamento}{Departamento de Ciencias e Ingeniería de la Computación}
\newcommand{\materia}{Organización de computadoras}
\newcommand{\cuatrimestre}{Segundo}

\newcommand{\RISC}{\textsc{risc}}
\newcommand{\CISC}{\textsc{cisc}}
\newcommand{\OCUNS}{\textsc{ocuns}}
\newcommand{\PC}{\textsf{PC}}
\newcommand{\SP}{\textsf{SP}}

% ENCABEZADO GENERAL
% DeptoLogo - Departamento - Universidad - Cuatrimestre - Año - UniLogo
\newcommand{\header}{
	\begin{center}
		\begin{tabular}{lcl}
		\multirow{4}{*}
		{\smash{$\vcenter{\hbox{\includegraphics[width=.78in]{DeptoLogo}}}$}} & {\Large\textsc{\materia}} &  
		\multirow{4}{*}
		{\smash{$\vcenter{\hbox{\includegraphics[width=.88in]{UniLogo}}}$}} \\
		         & \departamento & \\
		         & \universidad & \\
		         & \textbf{\cuatrimestre\ Cuatrimestre de \number\year} & \\
		\end{tabular}\\[2em]
	\end{center}
}

% ENCABEZADO EDITABLE%
% Uso: \Titulo{Primer texto negrita}{Segundo texto normal}
% header + Texto 1 + Texto 2
\newcommand{\Titulo}[2]{ 
	\begin{center}
		\header
		 \textbf{#1} \\ #2 \\[4pt]
	\end{center}
}

% ENCABEZADO PRÁCTICOS CON NUMEROS
% Uso: \TPnumero{Número de Práctico}{Título de Práctico}
% header + Trabajo Práctico Nº X + Titulo de Práctico
\newcommand{\TPnumero}[2]{ 
	\begin{center}
		\header
        Trabajo Práctico N$^{\circ}$ #1 \\ \textbf{#2} \\[4pt]
    \end{center}
}

% ENCABEZADO PRÁCTICOS CON LETRAS
% Uso: \TPletra{Letra de Práctico}{Título de Práctico}
% header + Trabajo Práctico 'X' + Titulo de Práctico
\newcommand{\TPletra}[2]{ 
	\begin{center}
		\header
		Trabajo Práctico #1 \\ \textbf{#2} \\[4pt]
	\end{center}
}

% ENCABEZADO PRÁCTICOS PROGRAMACIÓN CON NUMEROS
% Uso: \TPPnumero{Número de Práctico}{Título de Práctico}
% header + Trabajo Práctico de Programación Nº X + Titulo de Práctico
\newcommand{\TPPnumero}[2]{ 
	\begin{center}
		\header
		Trabajo Práctico de Programación N$^{\circ}$ #1 \\ \textbf{#2} \\[4pt]
	\end{center}
}

% ENCABEZADO PRÁCTICOS PROGRAMACIÓN CON LETRAS
% Uso: \TPPletra{Letra de Práctico}{Título de Práctico}
% header + Trabajo Práctico de Programación 'X' + Titulo de Práctico
\newcommand{\TPPletra}[2]{ 
	\begin{center}
		\header
		Trabajo Práctico de Programación #1 \\ \textbf{#2} \\[4pt]
	\end{center}
}

% ENCABEZADO PROYECTO CON NÚMERO
% Uso: \PRnumero{Número de Proyecto}{Título de Proyecto}
% header + Proyecto Nº X + Titulo de Proyecto
\newcommand{\PRnumero}[2]{
	\begin{center}
		\header
		Proyecto N$^{\circ}$ #1 \\ \textbf{#2} \\[4pt]
	\end{center}
}

\begin{document}

\Examen{Primer Examen Parcial}

\begin{center}
	\emph{Apague cualquier dispositivo electrónico en su poder y manténgalo guardado. No puede utilizar auriculares, ni calculadora. Lea todo el ejercicio antes de comenzar a desarrollarlo.}
\end{center}

\PEjercicio{1}Dado el número \textbf{decimal} $-298.5625$ llevar adelante los siguientes cambios de base:
\begin{enumerate}[a)]
	\item Convertirlo a \textbf{octal}, empleando el método de la \textbf{división} tanto para la parte entera como para la parte fraccionaria, expresando el resultado en \textbf{complemento a la base}, con 4 dígitos octales para la parte entera y 6 para la parte fraccionaria.

	\item Convertirlo a \textbf{binario} utilizando el método de la \textbf{multiplicación} tanto para la parte entera como para la parte fraccionaria, expresando el resultado en \textbf{complemento a la base disminuida}, con 12 bits para la parte entera y 6 bits para la parte fraccionaria.
\end{enumerate}

\PEjercicio{2}Considerando los números {\textbf{decimales}} $X = 1537$ e $Y = 2559$, llevar adelante las siguientes operaciones con una precisión de cuatro dígitos (incluido el signo), indicando claramente el resultado obtenido y la existencia o no de \emph{overflow}:
\begin{enumerate}[a)]
	\item Calcular $- X - Y$, trabajando en \textbf{hexadecimal} en \textbf{complemento a la base}.
	\item Calcular $ X + Y $, trabajando en \textbf{hexadecimal} en \textbf{complemento a la base disminuida}.
	\item Calcular $X - Y$, haciendo uso de un hardware que opera en una codificación \textbf{BCD Exceso-3} y \textbf{complemento a la base}, indicando claramente qué operación se está realizando en cada uno de los pasos intermedios.
\end{enumerate}

\PEjercicio{3} Considerando el Código Cíclico Redundante (CRC):

\begin{enumerate}[a)]
	\item Construir el mensaje $T(x)$ a transmitir asociado al mensaje de datos
	$M(x) = 110\,1011\,1011$ empleando el polinomio generador $G(x) = x^3 + x^2 + 1$.

	\item Suponiendo que durante la transmisión el mensaje $T(x)$ es modificado con un error $E(x)$ de tal forma que el receptor recibe el mensaje $T'(x) = 110\,0011\,0011\,0010$, determinar cómo opera el mecanismo de detección de errores y cuál es la conclusión que se alcanza.

	\item Comparando el mensaje transmitido $T(x)$ y el mensaje recibido $T'(x)$, ¿cuál es el desarrollo del polinomio de error $E(x)$? Sabiendo cuál fue el error exacto que existió, ¿cuál es la longitud de la ráfaga en error? y ¿a qué conclusión se puede arribar?
\end{enumerate}


\PEjercicio{4} Considerando el código Hamming mínima distancia 3 (Hamming básico), empleando paridad par y estando la secuencia ordenada de izquierda a derecha:

\begin{enumerate}[a)]

	\item Calcular los bits de código asociados al dato $0110\,1011$ y armar el codeword correspondiente que integra el dato y los bits calculados. 

	\item Considerando que el receptor recibe el codeword $1011\,1101\,1001$ que contiene los bits de dato y de código $C_i$. Recalcular los bits de código y determinar cuál es el síndrome.

	\item Determinar cómo trabaja el mecanismo de detección/corrección antes una política $d=2$, $c=1$ cuando al recalcular los bits de código de un codeword recibido se obtiene síndrome $1110$.

	\item Determinar cómo trabaja el mecanismo de detección/corrección antes una política $d=3$, $c=0$ cuando al recalcular los bits de código de un codeword recibido se obtiene síndrome $1010$.

	\item Determinar cómo trabaja el mecanismo de detección/corrección antes una política $d=3$, $c=0$ cuando al recalcular los bits de código de un codeword recibido se obtiene síndrome $0000$.

\end{enumerate}

\PEjercicio{5} Dadas las siguientes declaraciones de tipos:
\begin{verbatim}
	typedef void * tElemento;
	typedef struct nodo{
	   tElemento elemento;
	   struct nodo * padre;
	   struct nodo * hijo_izq;
	   struct nodo * hijo_der;
	} * tNodo;
	typedef struct abb{
	   unsigned int cant_elementos;
	   struct nodo * raiz;
	} * tABB;
\end{verbatim}
Implementar en lenguaje C, ACÁ DEBERÍAMOS SELECCIONAR ALGUNA DE LAS SIGUIENTES OPCIONES y la otra la podemos dejar para el RECU :) 
\begin{itemize}
	\item una función \texttt{void insertar(tABB a, tElemento e, int (*f)(void *,void *))}, que dado un \emph{Árbol Binario de Búsqueda} y una función de comparación \texttt{f}, inserte el elemento \texttt{e} en su correspondiente ubicación, haciendo un correcto uso de la memoria dinámica. Se considera que la función \texttt{f}, devuelve -1 si el orden del primer argumento es menor que el orden del segundo, 0 si el orden es el mismo, y 1 si el orden del primer argumento es mayor que el orden del segundo. \textbf{Deben hacer uso de malloc al crear un nuevo nodo, y moverse sobre la ED haciendo uso del comparador.}
\end{itemize}

\begin{itemize}
	\item una función \texttt{void eliminar\_raiz(tABB a, tElemento e, int (*f)(void *,void *))}, que dado un \emph{Árbol Binario de Búsqueda}, elimine el elemento almacenado en el nodo raíz del árbol, y reacomode la estructura tABB de forma tal que siga representando un árbol binario de búsqueda. \textbf{Deben hacer uso de free al eliminar un nodo hoja que pasa a ser nueva raiz, y moverse sobre la ED.}
\end{itemize}  

\end{document}

